\documentclass[12pt,a4paper]{article}

\usepackage[T1]{fontenc}
\usepackage[utf8]{inputenc}
\usepackage[polish]{babel}
\usepackage{lmodern}
\usepackage{geometry}
\usepackage{graphicx}
\usepackage{float}
\usepackage{pdflscape}
\usepackage{booktabs}
\usepackage{tabularx}
\usepackage{array}
\usepackage{xcolor}
\usepackage{amsmath}
\usepackage{listings}
\usepackage{enumitem}
\usepackage{hyperref}
\usepackage{caption}
\usepackage{microtype}

\geometry{margin=2.3cm}
\setlength{\parindent}{0pt}
\setlength{\parskip}{0.55em}
\setlength{\emergencystretch}{2em}
\setlist{nosep}
\hypersetup{
    colorlinks=true,
    linkcolor=blue!45!black,
    urlcolor=blue!55!black,
    pdftitle={Jednostka arytmetyczno-logiczna dla liczb w formacie zmiennoprzecinkowym},
    pdfauthor={[Imie i nazwisko]}
}

\pagestyle{plain}

\definecolor{codebg}{RGB}{246,247,249}
\definecolor{codeframe}{RGB}{210,214,220}
\definecolor{codecomment}{RGB}{50,120,70}
\definecolor{codekeyword}{RGB}{35,70,150}

\lstset{
    basicstyle=\ttfamily\small,
    backgroundcolor=\color{codebg},
    frame=single,
    rulecolor=\color{codeframe},
    keywordstyle=\color{codekeyword}\bfseries,
    commentstyle=\color{codecomment},
    showstringspaces=false,
    columns=fullflexible,
    keepspaces=true,
    breaklines=true,
    tabsize=2
}

\newcommand{\projekt}{Jednostka arytmetyczno-logiczna dla liczb w formacie zmiennoprzecinkowym}
\newcommand{\autorzy}{Filip Stępień, Bartłomiej Karkoszka, Rafał Grot}
\newcommand{\grupa}{1ID21A}

\newcommand{\zrzut}[4][6.5cm]{%
    \begin{figure}[H]
        \centering
        \IfFileExists{#2}{%
            \includegraphics[width=#3\textwidth,height=0.72\textheight,keepaspectratio]{#2}%
        }{%
            \fbox{%
                \begin{minipage}[c][#1][c]{0.9\textwidth}
                    \centering
                    \textbf{\Large MIEJSCE NA ZRZUT EKRANU}\\[0.8em]
                    \texttt{\detokenize{#2}}\\[0.8em]
                    #4
                \end{minipage}%
            }%
        }
        \caption{#4}
    \end{figure}
}

\newcommand{\zrzutpoziomy}[4][11cm]{%
    \begin{figure}[H]
        \centering
        \IfFileExists{#2}{%
            \includegraphics[width=#3\linewidth,height=#1,keepaspectratio]{#2}%
        }{%
            \fbox{%
                \begin{minipage}[c][#1][c]{0.92\linewidth}
                    \centering
                    \textbf{\Large MIEJSCE NA ZRZUT EKRANU}\\[0.8em]
                    \texttt{\detokenize{#2}}\\[0.8em]
                    #4
                \end{minipage}%
            }%
        }
        \caption{#4}
    \end{figure}
}

\begin{document}
\pagenumbering{gobble}

\begin{titlepage}
    \centering
    \vspace*{1.4cm}
    {\Large Politechnika Świętokrzyska\par}
    \vspace{0.35cm}
    {\large Wydział Elektrotechniki, Automatyki i Informatyki\par}
    \vspace{0.2cm}
    {\large Kierunek: Informatyka\par}

    \vfill
    {\Huge\bfseries \projekt\par}
    \vspace{0.8cm}
    {\Large Sprawozdanie z projektu\par}
    \vfill

    \begin{tabular}{>{\bfseries}rl}
        Autorzy: & \autorzy \\
        Grupa: & \grupa \\
        Przedmiot: & Projektowanie systemów wbudowanych \\
    \end{tabular}

    \vfill
    {\large Kielce, \the\year\par}
\end{titlepage}

\pagenumbering{roman}
\tableofcontents
\clearpage
\pagenumbering{arabic}

\section{Cel i zakres projektu}

Celem projektu było zaprojektowanie oraz przetestowanie w języku VHDL
jednostki arytmetyczno-logicznej wykonującej trzy operacje na liczbach
zapisanych w autorskim formacie zmiennoprzecinkowym:

\begin{itemize}
    \item dodawanie,
    \item odejmowanie,
    \item mnożenie.
\end{itemize}

Układ został zbudowany z niezależnych modułów, połączonych na schemacie
blokowym \texttt{demo.bdf}. Dwa liczniki generują okresowo zmieniające się
argumenty wejściowe. Jednostka ALU oblicza wynik wybranej operacji, a sygnał
\texttt{DONE} informuje o zakończeniu obliczeń. Taka organizacja pozwala
obserwować działanie projektu zarówno na przebiegach cyfrowych w ModelSimie,
jak i po przypisaniu wyjść do fizycznych wyprowadzeń układu FPGA.

Projekt obejmuje:

\begin{itemize}
    \item moduły obliczeniowe napisane w VHDL,
    \item schemat blokowy najwyższego poziomu wykonany w Quartusie,
    \item skrypt przygotowujący czytelne przebiegi w ModelSimie,
    \item automatyczny test porównujący wyniki sprzętowe z wartościami
          referencyjnymi.
\end{itemize}

\section{Środowisko i struktura plików}

Projekt przygotowano dla układu \textbf{Intel/Altera Cyclone IV E
EP4CE115F29C7}. Do syntezy wykorzystano Quartus II, natomiast symulację
funkcjonalną przeprowadzono w ModelSim-Altera.

Wszystkie źródła opisujące działanie układu znajdują się w katalogu
\texttt{src/}. Są to moduły VHDL odpowiedzialne za reprezentację liczb,
generowanie argumentów, sterowanie czasem oraz wykonanie operacji
arytmetycznych. Dzięki umieszczeniu kodu funkcjonalnego w jednym katalogu
można łatwo odróżnić go od plików projektu Quartusa, skryptów testowych
i wyników generowanych podczas kompilacji.

Plik \texttt{demo.bdf} jest graficznym schematem najwyższego poziomu.
Nie zawiera osobnego algorytmu obliczeniowego, lecz określa, jakie moduły
VHDL są użyte oraz w jaki sposób połączone są ich porty. Pliki
\texttt{*.bsf} przechowują symbole bloków widoczne w edytorze schematów
Quartusa. Definiują wygląd symbolu i rozmieszczenie jego wejść oraz wyjść,
natomiast właściwe działanie bloków pozostaje zapisane w źródłach VHDL.

\begin{table}[H]
    \centering
    \caption{Najważniejsze pliki projektu}
    \begin{tabularx}{\textwidth}{>{\raggedright\arraybackslash\ttfamily}p{0.38\textwidth}X}
        \toprule
        \normalfont Plik & Znaczenie \\
        \midrule
        demo.bdf & Graficzny schemat najwyższego poziomu i połączenia modułów. \\
        *.bsf & Symbole modułów VHDL używane w edytorze schematów Quartusa. \\
        fp\_alu.qpf, fp\_alu.qsf & Konfiguracja projektu Quartus. \\
        src/fp\_pkg.vhd & Definicja typu liczby i funkcji pomocniczych. \\
        src/fp\_counter.vhd & Generator kolejnych wartości argumentów. \\
        src/clk\_div.vhd & Generator impulsów o zadanej liczbie taktów na sekundę. \\
        src/start\_delay.vhd & Jednotaktowe opóźnienie rozpoczęcia operacji ALU. \\
        src/fp\_add\_sub.vhd & Sekwencyjne dodawanie i odejmowanie. \\
        src/fp\_mul.vhd & Sekwencyjne mnożenie. \\
        src/fp\_normalize.vhd & Normalizacja wyniku. \\
        src/alu.vhd & Wybór operacji i wspólny interfejs ALU. \\
        src/*\_bdf.vhd & Adaptery pomiędzy rekordami VHDL a magistralami BDF. \\
        test/wave\_scope.do & Konfiguracja i uruchomienie przebiegów. \\
        test/check\_counter\_\allowbreak{}results.do & Automatyczna kontrola wyników. \\
        \bottomrule
    \end{tabularx}
\end{table}

\section{Format danych}

\subsection{Reprezentacja liczby}

Argumenty oraz wynik są przechowywane jako rekord \texttt{fp\_t} złożony
z dwóch 16-bitowych liczb ze znakiem:

\begin{lstlisting}[language=VHDL,caption={Format liczby używany w projekcie}]
type fp_t is record
    mantissa : signed(15 downto 0);
    exponent : signed(15 downto 0);
end record;
\end{lstlisting}

Mantysa jest liczbą stałoprzecinkową w kodzie uzupełnień do dwóch,
w formacie Q1.15. Jeden bit reprezentuje znak i część całkowitą, a 15 bitów
część ułamkową. Wartość liczby oblicza się ze wzoru:

\begin{equation}
    x = \frac{M}{2^{15}} \cdot 2^E,
\end{equation}

gdzie \(M\) jest całkowitą wartością 16-bitowej mantysy, a \(E\) jest
wykładnikiem.

\begin{table}[H]
    \centering
    \caption{Przykładowe wartości formatu}
    \begin{tabular}{rrrr}
        \toprule
        Mantysa hex & \(M\) & Wykładnik & Wartość \\
        \midrule
        \texttt{1000} & 4096 & 0 & \(0{,}125\) \\
        \texttt{4000} & 16384 & 0 & \(0{,}5\) \\
        \texttt{E000} & -8192 & 0 & \(-0{,}25\) \\
        \texttt{4000} & 16384 & 1 & \(1{,}0\) \\
        \bottomrule
    \end{tabular}
\end{table}

Format jest prostszy od IEEE 754. Nie zawiera wartości specjalnych
\textit{NaN} i nieskończoności ani bitu niejawnego. Dzięki temu algorytmy
arytmetyczne są czytelne, a przebiegi łatwe do zinterpretowania.

\subsection{Normalizacja}

Niezerowy wynik jest normalizowany tak, aby wartość bezwzględna mantysy
znajdowała się w przedziale:

\begin{equation}
    0{,}5 \leq |\text{mantysa}| < 1.
\end{equation}

Jeżeli mantysa jest zbyt mała, moduł \texttt{fp\_normalize} przesuwa ją
w lewo i jednocześnie zmniejsza wykładnik. Gdy wynik jest równy zero,
zarówno mantysa, jak i wykładnik są zerowane.

\section{Architektura systemu}

\subsection{Interfejs najwyższego poziomu}

\begin{table}[H]
    \centering
    \caption{Porty schematu \texttt{demo.bdf}}
    \begin{tabularx}{\textwidth}{>{\raggedright\arraybackslash\ttfamily}p{0.31\textwidth}p{0.11\textwidth}X}
        \toprule
        \normalfont Port & Kierunek & Znaczenie \\
        \midrule
        CLOCK & wejście & Zegar systemowy 50 MHz. \\
        RESET & wejście & Asynchroniczny reset aktywny stanem wysokim. \\
        OP[1..0] & wejście & Wybór operacji: 00 -- dodawanie, 01 -- odejmowanie, 10 -- mnożenie. \\
        A\_MANTISSA[15..0] & wyjście & Mantysa pierwszego argumentu. \\
        A\_EXPONENT[15..0] & wyjście & Wykładnik pierwszego argumentu. \\
        B\_MANTISSA[15..0] & wyjście & Mantysa drugiego argumentu. \\
        B\_EXPONENT[15..0] & wyjście & Wykładnik drugiego argumentu. \\
        RESULT\_MANTISA[15..0] & wyjście & Mantysa wyniku. Nazwa portu jest zgodna ze schematem. \\
        RESULT\_EXPONENT[15..0] & wyjście & Wykładnik wyniku. \\
        DONE & wyjście & Jednotaktowy impuls zakończenia operacji. \\
        \bottomrule
    \end{tabularx}
\end{table}

\clearpage
\begin{landscape}
\subsection{Przepływ sygnałów}

Generator \texttt{clk\_div} tworzy impuls \texttt{tick}. Impuls zwiększa
licznik A. Po wykonaniu pełnego obiegu przez licznik A jego sygnał
\texttt{carry} zwiększa licznik B. W ten sposób układ automatycznie przechodzi
przez wszystkie pary argumentów. Ten sam impuls, opóźniony o jeden takt przez
\texttt{start\_delay}, rozpoczyna obliczenie ALU.

\zrzutpoziomy[11cm]{figures/quartus_blocks.png}{0.98}
{Schemat \texttt{demo.bdf} w edytorze Block Diagram/Schematic File programu Quartus. Widoczne są moduły projektu, magistrale danych oraz połączenia zegara, resetu, \texttt{tick}, \texttt{carry} i \texttt{start}.}
\end{landscape}
\clearpage

\section{Opis modułów}

\subsection{Generator impulsów \texttt{clk\_div}}

Moduł nie tworzy nowej domeny zegarowej. Zlicza zbocza narastające zegara
\texttt{CLOCK} i generuje jednotaktowy sygnał zezwolenia \texttt{tick}.
Liczba taktów pomiędzy impulsami wynosi:

\begin{equation}
    N = \frac{\texttt{CLOCK\_HZ}}{\texttt{TICKS\_PER\_SECOND}}.
\end{equation}

Na schemacie przyjęto:

\begin{align}
    \texttt{CLOCK\_HZ} &= 50\,000\,000,\\
    \texttt{TICKS\_PER\_SECOND} &= 5\,000\,000,\\
    N &= 10.
\end{align}

Przy zegarze 50 MHz okres zegara wynosi 20 ns, dlatego nowy impuls
\texttt{tick} pojawia się co:

\begin{equation}
    10 \cdot 20\ \text{ns} = 200\ \text{ns}.
\end{equation}

Parametr można zmniejszyć podczas pracy z fizycznym oscyloskopem, aby zmiany
były wolniejsze. Zwiększona wartość użyta w projekcie znacznie skraca czas
symulacji.

\subsection{Generatory argumentów \texttt{fp\_counter}}

Oba liczniki mają parametr \texttt{WIDTH=4}. Wykorzystują 16 stanów i
generują następujący cykl mantysy:

\begin{equation}
0,\ \frac18,\ \frac28,\ldots,\frac78,\ -1,\ -\frac78,\ldots,-\frac18.
\end{equation}

Wykładnik argumentów jest równy zero. Wartości mantysy zmieniają się zatem od
\(-1\) do \(0{,}875\) ze skokiem \(0{,}125\). Po przejściu z ostatniego stanu
do zera licznik A wystawia \texttt{carry}. Sygnał ten jest podłączony do
wejścia \texttt{en} licznika B.

Takie połączenie działa jak dwucyfrowy licznik:

\begin{itemize}
    \item A jest szybką cyfrą i zmienia się przy każdym impulsie
          \texttt{tick},
    \item B jest wolną cyfrą i zmienia się po pełnych 16 krokach A,
    \item w czasie jednego pełnego przebiegu sprawdzanych jest
          \(16 \cdot 16 = 256\) par argumentów.
\end{itemize}

\subsection{Opóźnienie startu \texttt{start\_delay}}

Licznik A aktualizuje argument na zboczu zegara wywołanym przez
\texttt{tick}. Jednostka ALU nie powinna w tym samym zboczu pobierać starej
wartości. Dlatego \texttt{start\_delay} opóźnia rozpoczęcie operacji o jeden
takt zegara:

\begin{center}
\begin{tabular}{c|cccc}
    Takt & \(k\) & \(k+1\) & \(k+2\) & \(k+3\) \\
    \hline
    \texttt{tick} & 1 & 0 & 0 & 0 \\
    zmiana A/B & tak & stabilne & stabilne & stabilne \\
    \texttt{start} & 0 & 1 & 0 & 0 \\
\end{tabular}
\end{center}

ALU rozpoczyna zatem pracę dopiero wtedy, gdy oba argumenty są już stabilne.

\subsection{Adaptery \texttt{fp\_counter\_bdf} i \texttt{alu\_bdf}}

Wewnętrzne moduły VHDL używają rekordów typu \texttt{fp\_t}. Edytor
schematów Quartusa wygodniej obsługuje osobne magistrale
\texttt{std\_logic\_vector}. Adaptery rozdzielają rekordy na mantysę
i wykładnik oraz składają je ponownie przed przekazaniem do ALU.

Adaptery nie wykonują obliczeń i nie wprowadzają opóźnienia. Stanowią jedynie
warstwę dopasowującą interfejs kodu VHDL do symboli użytych w pliku BDF.
Na ich podstawie Quartus korzysta z symboli
\texttt{fp\_counter\_bdf.bsf} i \texttt{alu\_bdf.bsf}, które można wstawić
na schemat oraz połączyć zwykłymi przewodami i magistralami. Pozostałe
symbole, między innymi \texttt{clk\_div.bsf} i
\texttt{start\_delay.bsf}, odpowiadają bezpośrednio portom modułów VHDL.

Rozdzielenie to pozwala zachować czytelny kod źródłowy z rekordem
\texttt{fp\_t}, a jednocześnie przedstawiać najważniejsze połączenia całego
systemu w graficznej postaci na schemacie \texttt{demo.bdf}.

\subsection{Jednostka nadrzędna \texttt{alu}}

Moduł \texttt{alu} uruchamia równolegle:

\begin{itemize}
    \item blok dodawania/odejmowania \texttt{fp\_add\_sub},
    \item blok mnożenia \texttt{fp\_mul}.
\end{itemize}

Sygnał \texttt{OP} wybiera wynik i sygnał \texttt{done} właściwego bloku.
Kodowanie operacji przedstawia tabela~\ref{tab:op}.

\begin{table}[H]
    \centering
    \caption{Kodowanie wejścia \texttt{OP}}
    \label{tab:op}
    \begin{tabular}{cl}
        \toprule
        \texttt{OP} & Operacja \\
        \midrule
        \texttt{00} & \(A+B\) \\
        \texttt{01} & \(A-B\) \\
        \texttt{10} & \(A\cdot B\) \\
        \texttt{11} & kod niewykorzystywany \\
        \bottomrule
    \end{tabular}
\end{table}

\subsection{Dodawanie i odejmowanie}

Blok \texttt{fp\_add\_sub} jest maszyną stanów. Operacja przebiega
w czterech etapach:

\begin{table}[H]
    \centering
    \caption{Stany bloku dodawania i odejmowania}
    \begin{tabularx}{\textwidth}{>{\ttfamily}p{0.2\textwidth}X}
        \toprule
        \normalfont Stan & Działanie \\
        \midrule
        S\_IDLE & Oczekiwanie na \texttt{start} i zapamiętanie argumentów. \\
        S\_ALIGN & Wyrównanie wykładników przez przesunięcie mantysy liczby o mniejszym wykładniku. \\
        S\_COMPUTE & Dodanie lub odjęcie rozszerzonych mantys oraz obsługa przepełnienia. \\
        S\_NORMALIZE & Normalizacja, zapis wyniku i wygenerowanie \texttt{done}. \\
        \bottomrule
    \end{tabularx}
\end{table}

Przed wykonaniem działania wykładniki muszą być równe. Mantysa liczby
o mniejszym wykładniku jest przesuwana arytmetycznie w prawo, a jej wykładnik
zwiększany. Po wyrównaniu wykonywane jest dodawanie lub odejmowanie na
17 bitach. Dodatkowy bit umożliwia wykrycie wyniku wykraczającego poza zakres
Q1.15. W takim przypadku mantysa jest przesuwana w prawo, a wykładnik
zwiększany o jeden.

\subsection{Mnożenie}

Blok \texttt{fp\_mul} działa w trzech stanach:

\begin{enumerate}
    \item zapamiętuje argumenty,
    \item mnoży dwie 16-bitowe mantysy i dodaje wykładniki,
    \item normalizuje wynik i zgłasza zakończenie.
\end{enumerate}

Iloczyn dwóch liczb Q1.15 ma 32 bity i skalę Q2.30. Aby wrócić do formatu
Q1.15, wybierane są bity \texttt{30 downto 15}. Wartość przed normalizacją
można zapisać jako:

\begin{equation}
    M_R = \frac{M_A \cdot M_B}{2^{15}},
    \qquad E_R = E_A + E_B.
\end{equation}

Po tej operacji wynik jest przekazywany do wspólnego modułu normalizującego.

\section{Przebieg działania}

Po zwolnieniu resetu system pracuje autonomicznie:

\begin{enumerate}
    \item \texttt{clk\_div} odmierza dziesięć taktów zegara,
    \item impuls \texttt{tick} zwiększa argument A,
    \item po pełnym obiegu A zwiększany jest argument B,
    \item takt później \texttt{start\_delay} generuje \texttt{start},
    \item ALU zapamiętuje A, B oraz wykonuje operację wybraną przez
          \texttt{OP},
    \item po zakończeniu wystawia wynik oraz jednotaktowy impuls
          \texttt{DONE}.
\end{enumerate}

Okres 200 ns pomiędzy kolejnymi impulsami \texttt{tick} jest większy od
latencji każdego bloku arytmetycznego, dlatego bieżąca operacja kończy się
przed rozpoczęciem następnej.

\section{Kompilacja w Quartusie}

\subsection{Procedura}

\begin{enumerate}
    \item Otworzyć plik projektu \texttt{fp\_alu.qpf}.
    \item Sprawdzić, czy jednostką najwyższego poziomu jest \texttt{demo}.
    \item Wybrać \texttt{Processing -> Start Compilation}.
    \item Po zakończeniu sprawdzić liczbę błędów w oknie komunikatów.
    \item Do symulacji funkcjonalnej wygenerować plik
          \texttt{simulation/qsim/fp\_alu.vo}.
\end{enumerate}

Pełna kompilacja sprawdzonej wersji projektu kończy się bez błędów.

\zrzut{figures/quartus_compilation.png}{0.92}
{Wynik kompilacji projektu w Quartusie. Kompilacja jednostki najwyższego poziomu \texttt{demo} zakończyła się bez błędów.}

\section{Symulacja przebiegów w ModelSimie}

\subsection{Uruchomienie}

Skrypt należy uruchomić z katalogu \texttt{simulation/qsim}:

\begin{lstlisting}[language=tcl,caption={Uruchomienie symulacji przebiegów}]
do ../../test/wave_scope.do
\end{lstlisting}

Skrypt automatycznie:

\begin{itemize}
    \item kompiluje netlistę \texttt{fp\_alu.vo},
    \item ładuje projekt \texttt{work.demo},
    \item generuje zegar 50 MHz,
    \item wymusza reset i kolejne wartości \texttt{OP},
    \item dodaje najważniejsze sygnały do okna Wave,
    \item ustawia mantysy jako analogowe przebiegi schodkowe,
    \item uruchamia symulację na 50 \(\mu\)s.
\end{itemize}

\begin{table}[H]
    \centering
    \caption{Wymuszenia użyte przez \texttt{wave\_scope.do}}
    \begin{tabular}{lll}
        \toprule
        Czas & Sygnał & Wartość \\
        \midrule
        0--200 ns & \texttt{RESET} & 1 \\
        od 200 ns & \texttt{RESET} & 0 \\
        0--15 \(\mu\)s & \texttt{OP} & 00 -- dodawanie \\
        15--30 \(\mu\)s & \texttt{OP} & 01 -- odejmowanie \\
        30--50 \(\mu\)s & \texttt{OP} & 10 -- mnożenie \\
        \bottomrule
    \end{tabular}
\end{table}

\clearpage
\begin{landscape}
\subsection{Interpretacja przebiegów}

Na wykresie A tworzy szybką piłę schodkową, ponieważ jego licznik jest
zwiększany przy każdym impulsie \texttt{tick}. B zmienia się 16 razy wolniej,
gdyż jest zwiększane dopiero po pełnym cyklu A. Wynik zależy od operacji
wskazanej przez \texttt{OP}:

\begin{itemize}
    \item dla \texttt{OP=00} wynik jest sumą A i B; ponieważ B pozostaje
          chwilowo stałe, przebieg przypomina A przesunięte o wartość B,
    \item dla \texttt{OP=01} wynik jest różnicą A i B; odejmowanie stałego
          fragmentu B przesuwa przebieg A w dół,
    \item dla \texttt{OP=10} B skaluje przebieg A; dodatnie B zachowuje jego
          kierunek, ujemne odwraca znak, a wartość bezwzględna B określa
          amplitudę.
\end{itemize}

Mantysę i wykładnik należy analizować łącznie. Sama mantysa może zmaleć po
normalizacji, mimo że rzeczywista wartość wyniku wzrosła. Przykładowo liczba
1 jest reprezentowana jako mantysa \(0{,}5\) i wykładnik 1. Z tego powodu
widoczny kształt mantysy wyniku może zawierać dodatkowe uskoki w chwilach
zmiany wykładnika, mimo że rzeczywista wartość operacji zmienia się zgodnie
z opisanymi zależnościami. Schodkowy wygląd wynika natomiast z tego, że
ModelSim rysuje kolejne cyfrowe wartości magistrali jako przebieg
\textit{Analog (Step)}.

\zrzutpoziomy[8.2cm]{figures/modelsim_waveforms.png}{0.98}
{Okno Wave po wykonaniu \texttt{wave\_scope.do}. Mantysy A, B i wyniku przedstawiono jako przebiegi \textit{Analog (Step)}; widoczne są również wykładniki, \texttt{DONE}, \texttt{OP}, \texttt{RESET} i \texttt{CLOCK}.}
\end{landscape}
\clearpage

\section{Automatyczna weryfikacja obliczeń}

\subsection{Sposób działania testu}

Samo obejrzenie przebiegów pozwala ocenić ogólny kształt sygnałów, ale nie
gwarantuje poprawności każdej wartości. Dlatego przygotowano skrypt
\texttt{check\_counter\_results.do}.

Uruchomienie:

\begin{lstlisting}[language=tcl,caption={Uruchomienie automatycznego testu}]
do ../../test/check_counter_results.do
\end{lstlisting}

Po każdym impulsie \texttt{DONE} skrypt:

\begin{enumerate}
    \item odczytuje mantysy i wykładniki A, B oraz wyniku,
    \item zamienia zapis sprzętowy na wartości rzeczywiste,
    \item oblicza oczekiwany wynik dla aktualnego \texttt{OP},
    \item porównuje wynik sprzętowy z referencyjnym,
    \item zapisuje wiersz \texttt{PASS} albo \texttt{FAIL}.
\end{enumerate}

Wartość odczytana z portów jest rekonstruowana zgodnie ze wzorem:

\begin{equation}
    x_{\text{ModelSim}} =
    \frac{\operatorname{signed}(\text{mantissa})}{32768}
    \cdot 2^{\operatorname{signed}(\text{exponent})}.
\end{equation}

Porównanie wykorzystuje tolerancję \(10^{-9}\). Dla wartości występujących
w tym teście nie jest potrzebne dodatkowe zaokrąglanie, ponieważ są one
wielokrotnościami potęg liczby 2 i mają dokładną reprezentację binarną.

\subsection{Przykładowe przypadki}

\begin{table}[H]
    \centering
    \caption{Przykładowe wyniki sprawdzane przez skrypt}
    \begin{tabular}{lrrrr}
        \toprule
        Operacja & A & B & Oczekiwano & Otrzymano \\
        \midrule
        ADD & \(0{,}875\) & \(0{,}125\) & \(1{,}000\) & \(1{,}000\) \\
        SUB & \(-0{,}750\) & \(0{,}500\) & \(-1{,}250\) & \(-1{,}250\) \\
        MUL & \(0{,}625\) & \(-0{,}875\) & \(-0{,}546875\) & \(-0{,}546875\) \\
        \bottomrule
    \end{tabular}
\end{table}

Pełny dziennik jest zapisywany do pliku:

\begin{center}
    \texttt{build/counter\_results.log}
\end{center}

W zweryfikowanej wersji projektu otrzymano:

\begin{lstlisting}[caption={Podsumowanie automatycznego testu}]
SUMMARY: checked=248 passed=248 failed=0
COUNTER/ALU CHECK PASSED
\end{lstlisting}

Liczba 248 odpowiada operacjom zakończonym w czasie 50 \(\mu\)s przy
zastosowanych wymuszeniach. Nie jest to ograniczenie liczby kombinacji
wejściowych, lecz rezultat wybranego czasu symulacji.

\zrzut{figures/modelsim_checker.png}{0.94}
{Okno Transcript po wykonaniu \texttt{check\_counter\_results.do}. Końcowe podsumowanie zawiera \texttt{checked=248 passed=248 failed=0} oraz komunikat \texttt{COUNTER/ALU CHECK PASSED}.}

\section{Uzyskane rezultaty}

Przeprowadzona kompilacja i symulacja potwierdziły:

\begin{itemize}
    \item poprawne generowanie argumentów A i B,
    \item właściwe opóźnienie sygnału startu względem zmiany argumentów,
    \item poprawne działanie dodawania, odejmowania i mnożenia,
    \item prawidłową normalizację mantysy i aktualizację wykładnika,
    \item generowanie impulsu \texttt{DONE} po zakończeniu obliczenia,
    \item zgodność 248 kolejnych wyników z modelem referencyjnym.
\end{itemize}

Przebiegi analogowe w ModelSimie ułatwiają ocenę działania systemu:
argument A jest szybką piłą, B wolną piłą, a wynik tworzy kształt zależny
od operacji. Automatyczny test uzupełnia tę ocenę o dokładne sprawdzenie
wartości liczbowych.

\section{Ograniczenia i możliwy rozwój}

Aktualna wersja jest demonstratorem jednostki ALU, dlatego ma kilka świadomie
przyjętych ograniczeń:

\begin{itemize}
    \item format nie jest zgodny z IEEE 754,
    \item nie występują wartości specjalne ani obsługa wyjątków,
    \item wykładniki argumentów z liczników są stale równe zero,
    \item nie zastosowano bitów ochronnych i rozbudowanego zaokrąglania,
    \item kod \texttt{OP=11} nie realizuje operacji,
    \item przed implementacją na płytce należy przypisać piny i dodać
          ograniczenia czasowe.
\end{itemize}

Projekt można rozwinąć przez dodanie dzielenia, testów dla różnych
wykładników, obsługi przepełnienia wykładnika, trybów zaokrąglania oraz
interfejsu umożliwiającego podawanie argumentów z zewnątrz.

\section{Skrócona instrukcja odtworzenia wyników}

\begin{enumerate}
    \item Otworzyć \texttt{fp\_alu.qpf} w Quartusie.
    \item Uruchomić pełną kompilację projektu.
    \item Upewnić się, że w \texttt{simulation/qsim} znajduje się
          \texttt{fp\_alu.vo}.
    \item Uruchomić ModelSim-Altera w katalogu \texttt{simulation/qsim}.
    \item Wpisać \texttt{do ../../test/wave\_scope.do}, aby obejrzeć
          przebiegi.
    \item Wpisać \texttt{do ../../test/check\_counter\_results.do}, aby
          wykonać kontrolę liczbową.
    \item Sprawdzić podsumowanie w oknie Transcript oraz dziennik
          \path{build/counter_results.log}.
\end{enumerate}

\section{Podsumowanie}

W ramach projektu wykonano modułową jednostkę ALU operującą na liczbach
z 16-bitową mantysą Q1.15 i 16-bitowym wykładnikiem. Układ realizuje
dodawanie, odejmowanie oraz mnożenie, a generatory argumentów pozwalają na
ciągłą obserwację jego działania bez zewnętrznego źródła danych.

Połączenie schematu blokowego Quartusa, czytelnych przebiegów ModelSim oraz
automatycznej kontroli wartości zapewnia trzy poziomy weryfikacji:
strukturalny, wizualny i liczbowy. Wszystkie sprawdzone przypadki zakończyły
się wynikiem poprawnym.

\end{document}
